\bookStart{Muspell}[Muspilli]
\setBookCode{Muspilli}

\begin{flushright}%
\textbf{Dating:} C9th

\textbf{Meter:} \Fornyrdislag%para
\end{flushright}%

\section{Introduction}

The \textbf{Muspell} (\Muspilli) is an Old High German Christian poem dealing with the Day of Judgment.

\subsection{Preservation}

\Muspilli\ survives in a single copy, found scribbled in a Latin-language theological manuscript from the 820s \ce\ with signum CLM 14098 (within the confines of the pres. poem denoted as \textbf{A}); since the poem is marginalia, the dating of the manuscript can unfortunately only serve as a \emph{terminus post quem} for the poem itself.

\subsection{Content}

\Muspilli\ is clearly a Catholic polemic, describing the events of the Judgment Day and advising its readers to do penance and give alms to assure their safety.  The poem is hardly a poetic masterpiece and comes off as quite stilted even in the original, although its meter cannot be called irregular when compared to the other High German alliterative poems \textlink{Hildebrandslied} and \textlink{Wessobrunn}.  In the English translation there is no poetic quality left.

The use of occasional end rhyme (see note to ll. 60--61) suggests a relation to Otfrid’s \emph{Evangelienbuch} (written 863--871 \ce), as does the exact correspondence between \Muspilli\ 14 and \emph{Evangelienbuch} 1.18.9.  Whatever the direction of influence, the author of \Muspilli\ surely belonged to the same monastic C9th milieu as Otfrid.

\subsubsection{Language}

The dialect is that of the southern High German area, as seen by the consistent application of the most extensive form of the second sound shift, where \emph{g, b, k} change to \emph{k, p, ch}.  That this was the case at the time of composition is proven by the fact that Germanic roots originally beginning with \emph{b} consistently alliterate with Latin borrowings beginning with \emph{p}, namely in:

\begin{itemize}
  \item l. 16: \emph{pú} (= OS \emph{bú}) : \emph{pardísu} (< Latin \emph{paradīsum}),
  \item l. 21: \emph{piutit} (= OS \emph{biudid}) : \emph{pehhes} (< Latin \emph{pix}), \emph{pína} (< Latin \emph{poena}),
  \item l. 25: \emph{prinnan} (= OS \emph{brinnan}), \emph{palw-} (= OS \emph{balu}) : \emph{pehhe}.
\end{itemize}

Interestingly, the alliteration also shows that the poet retained old \emph{h-} before \emph{l} (l. 72), by extension almost certainly also before \emph{r} and \emph{n}, and probably also before \emph{w} (l. 7).  This sound is, however, consistently omitted by the scribe.

Due to the low quality of the digitally available scans of the already very poorly preserved manuscript, the pres. ed. has had to rely partly on the text found in the 17th edition of Braune’s \emph{Althochdeutsche Lesebuch} published in 1994.%TODO: add when true:  All critical choices have, however, been controllied by the editor against scans of the manuscript.

\sectionline
\newpage

\section{Text}

\bvg\bva[]%
{[...]} sín \alst{t}ak pi·kweme, \hld\ daz er \alst{t}ouwan skal. &
\emph{H}wanta \alst{s}ár só sih diu \alst{s}êla \hld\ in den \alst{s}ind ar·hęvit, &
ęnti sí den \alst{l}íh-hamun \hld\ \edtext{\alst{l}ikkan lázzit}{\Bfootnote{The technical double alliteration in the second half-line is unusual, but probably not due to any scribal corruption.}}, &
só kwimit ęin \alst{h}ęri \hld\ fona \alst{h}imil-zungalon; &
daz andar fona \alst{p}ehhe: \hld\ dár \alst{p}ágant siu umpi.\eva

\bvb {[...]} his day will come when he must die. \\
For as soon as the soul lifts itself on its way \\
and leaves the body behind to lie, \\
then comes an army from the heavenly stars, \\
the other from the Pitch; there they fight over it.\evb\evg


\bvg\bva[][6]%
\alst{S}orgén mak diu \alst{s}êla, \hld\ unzi diu \alst{s}uona ar·gét, &
za \edtext{\emph{\alst{h}}wederemo}{\Afootnote{\emph{wederemo} \textbf{A}}\Bfootnote{Restoration of the initial \emph{h-} is not strictly required for the line to alliterate properly, but is done on the basis of l. 72.}} \alst{h}ęrje \hld\ si gi·\alst{h}alót werde. &
\emph{H}wanta ipu sia daz \alst{S}atanazses \hld\ ki·\alst{s}indi ki·winnit, &
daz \alst{l}ęitit sia sár \hld\ dár iru \alst{l}ęid wirdit, &
in \alst{f}uir ęnti in \alst{f}instrí: \hld\ daz ist rehto \alst{v}irin-líh ding.\eva

\bvb The soul must worry until the time of atonement \\
about which of the armies are to receive it, \\
for if Satan’s following wins it \\
then it is soon led whither it must face torment \\
into fire and into darkness.  That is truly a wicked thing.\evb\evg


\bvg\bva[][11]%
Upi sia \edtext{*avar}{\Afootnote{\emph{hauar} \textbf{A}}} ki·\alst{h}alónt die \hld\ die dár fona \alst{h}imile kwemant, &
ęnti sí dero \alst{ę}ngilo \hld\ \alst{ęi}gan wirdit, &
\edtext{die pringent sia sár úf \hld\ in himilo ríhi:}{\Bfootnote{This line lacks alliteration, but no obvious emendation presents itself.}} &
\edtext{dár ist \alst{l}íp áno tôd, \hld\ \edtext{\alst{l}i\emph{oh}t}{\Afootnote{\emph{lihot} \textbf{A}}} áno \edtext{finst\emph{r}í}{\Afootnote{\emph{finsti} \textbf{A}}}}{\Bfootnote{This line also appears in Otfrid’s \emph{Evangelienbuch} 1.18.9, in the form: \emph{Thár ist líb ána tôd, \hld\ lioht ána finstri.}  It is one of Otfrid’s rhymeless lines where alliteration compensates for the expected end-rhyme.  For the relevance of this shared line to the relation between \Muspilli\ and \emph{Evangelienbuch} see Introduction above.}}, &
\alst{s}ęlida áno \edtext{\alst{s}orgu\emph{n: \hld\ dár n·ist}}{\Afootnote{emend.; partly illegible in \textbf{A}}} neo-man \alst{s}iuh. &
Denne der man in \edtext{\alst{p}ar\emph{dí}su}{\Afootnote{partly illegible in \textbf{A}}} \hld\ \alst{p}ú ki·winnit, &
\alst{h}ús in \alst{h}imile, \hld\ dár kwimit imo \alst{h}ilfa ki·nuok.\eva

\bvb If, however, those who come from heaven receive it \\
and the angels are to have it for themselves, \\
they soon bring it up into the realm of the Heavens. \\
There life is without death, light without darkness, \\
happiness without sorrows; there no man is sick. \\
When the man wins his dwelling in Paradise, \\
his house in heaven, he will get help aplenty.\evb\evg


\bvg\bva[][18]%
Pi·diu ist durft \alst{m}ihhil allero \alst{m}anno \edtext{\emph{h}we-líhemo}{\Afootnote{end 61r; the text picks back up at 119v.}}, \hld\ daz in es sín \alst{m}uot ki·spane, &
daz er \alst{k}otes willun \hld\ \alst{k}erno tu\emph{e} &
ęnti \alst{h}ęlla fuir \hld\ \alst{h}arto wíse, &
\alst{p}ehhes \alst{p}ína: \hld\ dár \alst{p}iutit der Satanasz altist &
\alst{h}ęizzan lauk. \hld\ Só mak \alst{h}ukkan za diu, &%NOTE: za] 120r
\alst{s}orgén dráto, \hld\ der sih \alst{s}untigen węiz.\eva

\bvb Therefore it is of great need for every man that he persuade his heart \\
that it may gladly do the will of God \\
and stalwartly avoid the fire of Hell, \\
the torture of the Pitch; there Satan always offers \\
hot flame.  So he must think about that, \\
worry greatly, whoever knows that he sins.\evb\evg


\bvg\bva[][24]%
Wê demo in \alst{v}instrí skal \hld\ síno \alst{v}iriná stúén, &
\alst{p}rinnan in \edtext{\alst{p}\emph{e}hhe}{\Afootnote{\emph{phhe} \textbf{A}}}: \hld\ daz ist rehto \alst{p}alwík dink, &
daz der man \alst{h}arét ze gote \hld\ ęnti imo \alst{h}ilfa ni kwimit. &
\alst{W}ánit sih ki·náda \hld\ diu \alst{w}ênaga sêla: &%NOTE: wánit] 120v
ni ist in ki·\alst{h}uktin \hld\ \alst{h}imiliskin gote, &
\emph{h}wanta hiar in \alst{w}er-olti \hld\ after ni \alst{w}erkóta.\eva

\bvb Woe to him who in the darkness shall atone for his wickedness, \\
burn in the Pitch; that is truly an accursed thing, \\
that the man calls out to God and no help for him comes. \\
It expects mercy, that wretched soul; \\
it is not in the thoughts of the heavenly God, \\
for here in this world it left no works behind.\evb\evg


\bvg\bva[][30]%
Só denne der \alst{m}ahtigo khunink \hld\ daz \alst{m}ahal ki·pannit, &
dara skal \alst{k}weman \hld\ \alst{kh}unno ki·líhaz: &
denne ni ki·tar \alst{p}arno nohhęin \hld\ den \alst{p}an furi·sizzan, &
ni allero \alst{m}anno \emph{h}we-líh; \hld\ ze demo \alst{m}ahale skuli. &
Dár skal er vora demo \alst{r}íhhe \hld\ az \alst{r}ahhu stantan, &
pí daz er in \alst{w}er-olti eo \hld\ ki·\alst{w}erkót hapéta.\eva

\bvb So, when the mighty King summons the court \\
whither every race shall come, \\
then none of the children dare ignore the summons \\
nor whatsoever man; he must go to the court. \\
There shall he before that power stand on trial \\
for the works he has done earlier in the world.\evb\evg


\bvg\bva[][36]%
Daz hôrt’ ih \alst{r}ahhón \hld\ dia wer-olt-\alst{r}eht-wíson, &
daz skuli der \alst{A}nti-khristo \hld\ mit \alst{E}líase págan. &
Der \alst{w}arkh ist ki·\alst{w}áfanit, \hld\ denne wirdit untar in \alst{w}ík ar·hapan. &
\alst{Kh}ęnfun sint só \alst{k}ręftík; \hld\ diu \alst{k}ósa ist só mihhil. &
\alst{E}lías strítit \hld\ pí den \alst{ê}wigon líp, &
wili dén \alst{r}eht-kernón \hld\ daz \alst{r}íhhi ki·starkan: &
pi·diu skal imo \alst{h}elfan \hld\ der \alst{h}imiles ki·waltit.\eva

\bvb This I heard the righteous men of the world tell, \\
that the Anti-Christ shall fight with Elijah; \\
the outlaw is armed; then he will commence the battle. \\
The champions are so strong, the strife is so great; \\
Elijah fights for the eternal life, \\
wishes to strengthen the power of the righteous; \\
therefore he will have help from Him Who rules Heaven.\evb\evg


\bvg\bva[][43]%
Der \alst{A}nti-khristo \hld\ stét pí demo \alst{a}lt-fíante, &
stét pí demo \alst{S}atanase, \hld\ der inan var·\alst{s}enkan skal: &
pi·diu skal er in deru \alst{w}ík-stęti \hld\ \alst{w}unt pi·vallan &
ęnti in demo \alst{s}inde \hld\ \alst{s}iga-lôs werdan.\eva

\bvb The Anti-Christ stands beside the Ancient Fiend, \\
stands beside Satan who shall overcome him. \\
Therefore shall he on that battlefield fall wounded \\
and in the incursion be victoryless.\evb\evg


\bvg\bva[][47]%
Doh wánit des vila got-manno, &
daz Elías in demo \alst{w}íge \hld\ ar·\alst{w}artit werde, &%NOTE: werde] 121r
só daz \alst{E}líases pluot \hld\ in \alst{e}rda ki·triufit.\eva

\bvb Yet many Godly men expect \\
that Elijah in that battle will be maimed, \\
so that Elijah’s blood is driven into the earth.\evb\evg


\bvg\bva[][50]%
Só in·\alst{p}rinnant die \edtext{\alst{p}erga, \hld\ \alst{p}oum}{\lemma{perga \dots\ poum ‘mountains \dots\ tree’}\Bfootnote{Formulaic word-pair; see note to \textlink{Wessobrunn} 3.}} ni ki·stęntit &
\alst{ê}nihk in \alst{e}rdu, \hld\ \alst{a}há ar·truknént, &
muor var·\alst{s}wilhit sih, \hld\ \alst{s}wilizót lougiu der himil, &
\alst{m}áno vallit, \hld\ prinnit \alst{m}ittila-gart, &
\edtext{\alst{st}ên ni ki·\alst{st}ęntit*}{\Afootnote{add. \emph{ênikh in erdu} \textbf{A} is unmetrical dittography from l. 50b--51a.}}, \hld\ vęrit denne \alst{st}úa-tago in lant, &
\alst{v}ęrit mit diu \alst{v}uiru \hld\ \alst{v}iriho wísón.\eva

\bvb So the mountains burn, no tree stands \\
anywhere upon the earth, the rivers turn dry, \\
the swamp is swallowed up, the heaven glows with flame, \\
the moon falls, Middenyard burns, \\
no stone stands.  Then the Judgment Day comes upon the land, \\
comes with the fire to afflict the folk.\evb\evg


\bvg\bva[][56]%
Dár ni mak denae \alst{m}ák andremo \hld\ helfan vora demo \edtrans{\alst{M}úspille}{the Muspell}{\Bfootnote{Most likely ‘the Destruction’, the word which has given the poem its modern name.  It appears to be of Germanic extraction and corresponds to ON \emph{Muspell}, OS \emph{Múd-spell}, but the exact etymology is quite unclear; it seems to be a cpd. consisting of an unknown first element probably meaning ‘the world’ or ‘the earth’ + a derivative of the verb \emph{*spilþijaną} ‘to spoil, ruin; kill, murder’ (OE \emph{spillan, spildan}, ON \emph{spilla}, OS \emph{spildjan}).}}, &
denne daz \alst{p}ręita wasal \hld\ allaz var·\alst{p}rinnit, &
ęnti \alst{v}uir ęnti luft \hld\ iz allaz ar·\alst{f}urpit. &
\edtrans{\emph{H}wár ist denne diu \alst{m}arha, \hld\ dár man dár eo mit sínén \alst{m}ágon piehk?}{Where is then the ground where man always with his kinsmen fought?}{\Bfootnote{An example of the \emph{ubi sunt} motif common in mediæval literature from both Europe and the Arab world.  This literary motif, originally drawing on the Book of Baruch (3:16--19) and in Old Germanic poetry most famously found in \emph{The Wanderer} (ll. 92--96), illustrates the vanity of life by listing asking ‘where are they?’  The German poet here extends the motif to the Day of Judgment and takes the opportunity to comment on his own time by condemning the internecine feuds that plagued all mediæval monarchies.}} &
\edtext{Diu marha ist far·prunnan, \hld\ diu sêla stét pi·d\emph{w}ungan, &
ni węiz mit \emph{h}wiu puaze: \hld\ só vęrit sí za wíze.}{\lemma{Diu \dots\ wíze}\Bfootnote{In these two lines the poet replaces the usual alliteration with end-rhyme within each half-lines pair (\emph{prunnan} : \emph{dwungan --- puaze} : \emph{wíze}).  The very same meter, including the looseness of the rhymes, is used by Otfrid throughout the whole of his \emph{Evangelienbuch}, written some time between 863 and 871 \ce.  The direction of influence between \Muspilli\ and that work is uncertain owing to the difficulties of dating the present poem, for which see introduction above.}}\eva

\bvb There may no kinsman help another one before the Muspell \\
when the broad earth all burns up \\
and fire and air cleanse it all. \\
Where is then the ground where man always with his kinsmen fought? \\
The ground is burnt up; the soul is defeated, \\
knows no remedy; so it goes to punishment.\evb\evg


\bvg\bva[][62]%
Pi·diu ist demo \alst{m}anne só guot, \hld\ denne ’r ze demo \alst{m}ahale kwimit, &
daz er \alst{r}ahóno \emph{h}we-líha \hld\ \alst{r}ehto ar·tęile; &
denne ni darf er \alst{s}orgén, \hld\ denne er ze deru \alst{s}uonu kwimit. &
Ni \alst{w}ęiz der \alst{w}ênago man, \hld\ \emph{h}wie-líhan \alst{w}artil er habét, &
denne ’r mit den \alst{m}iatón \hld\ \alst{m}arrit daz rehta, &
daz der \alst{t}iuval dár pí \hld\ ki·\alst{t}arnit stęntit.\eva

\bvb Therefore it is so good for the man when he comes to the court \\
that he rightly consider every charge, \\
then he need not worry when he comes to the atonement. \\
The wretched man knows not what sort of overseer he has, \\
when with bribes he obstructs what is right, \\
that the Devil stands hidden by his side.\evb\evg


\bvg\bva[][68]%
\edtrans{Der}{\emph{He}}{\Bfootnote{The Devil, as just mentioned.}} hapét in \alst{r}uovu \hld\ \alst{r}ahóno \emph{h}we-líha, &
daz der man \alst{ê}r ęnti síd \hld\ \alst{u}piles ki·frumita, &
daz er iz allaz ki·\alst{s}agét, \hld\ denne er ze deru \alst{s}uonu kwimit; &
ni skolta síd \alst{m}anno nohhęin \hld\ \alst{m}iatun int·fáhan.\eva%NOTE: ni: 121v

\bvb \emph{He} will in silence keep every charge, \\
which the man, early or late, has committed, \\
so that he tells it all when he comes to the atonement; \\
for that sake should no man take bribes.\evb\evg


\bvg\bva[][72]%
Só daz \alst{h}imiliska \alst{h}orn \hld\ \edtrans{ki·\emph{\alst{h}}lútit}{sounded}{\Afootnote{\emph{kilutit} \textbf{A}}\Bfootnote{Restoration of the cluster \emph{hl-} is required by the alliteration; cf. l. 7.}} wirdit, &
ęnti sih der \alst{s}uanari \hld\ ana den \alst{s}ind ar·hęvit &
der dár suannan skal \hld\ tôten ęnti lepentén; &
denne \alst{h}ęvit sih mit imo \hld\ \alst{h}ęrjo męista, &
daz ist allaz só \alst{p}ald, \hld\ daz imo nio-man ki·\alst{p}ágan ni mak. &
Denne vęrit er ze deru \alst{m}ahal-stęti, \hld\ deru dár ki·\alst{m}arkhót ist: &
dár wirdit diu \alst{s}uona, \hld\ dia man dár io \alst{s}agéta.\eva

\bvb So the heavenly horn is sounded, \\
and the Judge lifts Himself on His way, \\
He who shall judge the dead and the living. \\
Then comes with Him the greatest of armies; \\
it is all so bold that no man may fight it. \\
Then He journeys to the place of the court which is there marked out; \\
there begins the atonement as was once foretold.\evb\evg


\bvg\bva[][79]%
Denne varant \alst{ę}ngila \hld\ \alst{u}per dio marha, &
\alst{w}ękhant deota, \hld\ \alst{w}íssant ze dinge. &
Denne skal \alst{m}anno gi·líh \hld\ fona deru \alst{m}oltu ar·stén, &
\alst{l}ôssan sih ar dero \alst{l}éwo vazzón: \hld\ \edtext{skal imo avar sín \alst{l}íp pi·kweman}{\lemma{skal imo \dots\ sín líp pi·kweman ‘he shall come to life’}\Bfootnote{Lit. ‘his life shall come to him.’}}, &
daz er sín \alst{r}eht allaz \hld\ ki·\alst{r}ahhón muozzi, &
ęnti imo after sínén \alst{t}átin \hld\ ar·\alst{t}ęilit werde.\eva

\bvb Then the angels journey over the ground, \\
they wake the nations, show them to the Thing. \\
Then shall every man arise out of the earth, \\
free himself from the burdens of the grave; further he shall come to life \\
that he may fully plead his case \\
and according to his deeds be considered.\evb\evg


\bvg\bva[][85]%
Denne der gi·\alst{s}izzit, \hld\ der dár \alst{s}uonnan skal &
ęnti ar·\alst{t}ęillan skal \hld\ \alst{t}ôtén ęnti kwekkhén, &
denne stét dár \alst{u}mpi \hld\ \alst{ę}ngilo męnigí, &
\alst{g}uotero \alst{g}omóno: \hld\ \alst{g}art ist só mihhil: &
dara kwimit ze deru \alst{r}ihtungu só vilo \hld\ dia dár ar \alst{r}ęstí ar·stént.\eva

\bvb Then sits he who there shall judge \\
and shall consider the dead and the living. \\
Then there around him stands a throng of angels, \\
{[and]} of good men.  The enclosure is so great; \\
there to that tribunal come so many who have arisen from their rest.\evb\evg


\bvg\bva[][90]%
Só dár \alst{m}anno nohhęin \hld\ wiht pi·\alst{m}ídan ni mak; &
dár skal denne \alst{h}ant sprehhan, \hld\ \alst{h}oupit sagén, &
allero \alst{l}ido \emph{h}we-líhk \hld\ unzi in den \alst{l}uzígun vinger, &
\emph{h}waz er untar desen \alst{m}annun \hld\ \alst{m}ordes ki·frumita.\eva

\bvb So there may no man avoid the smallest thing; \\
there the hand shall speak, the head shall tell, \\
{[and]} every single limb down to the little finger, \\
what sort of murder it committed in the midst these men.\evb\evg


\bvg\bva[][94]%
Dár ni ist eo só \alst{l}istík man \hld\ der dár io·wiht ar·\alst{l}iugan męgi, &
daz er ki·\alst{t}arnan męgi \hld\ \alst{t}áto dehhęina, &
niz al fora demo \alst{kh}uninge \hld\ ki·\alst{kh}undit werde, &
\alst{ú}zzan er \alst{i}z \hld\ mit \alst{a}lamusanu furi·męgi &
ęnti mit \alst{f}astún \hld\ dio \alst{v}iriná ki·puazti. &
Denne der \alst{p}aldét \hld\ der gi·\alst{p}uazzit hapét, &
denne ’r ze deru suonu kwimit.\eva

\bvb There no man is ever so clever that he may lie about the smallest thing, \\
that he may hide away any deed, \\
so that it would not fully be made known to the King, \\
unless he overcame it through the giving of alms \\
and through fasting did penance for the wickedness. \\
Then is he emboldened who has done penance, \\
when he comes to the atonement.\evb\evg


\bvg\bva[][101]%
Wirdit denne \alst{f}uri ki·tragan \hld\ daz \alst{f}rôno khrúki, &
dár der \alst{h}êligo Khrist \hld\ ana ar·\alst{h}angan ward. &
Denne augit er dio \alst{m}ásún, \hld\ dio er in deru \alst{m}ęnniskí an·fénk, &
dio er duruh desse \alst{m}an-kunnes \hld\ \alst{m}inna far·doléta.\eva

\bvb Then is brought forth the Lordly Cross, \\
whereon the Holy Christ was made to hang. \\
Then He makes seen the wounds which He as a man received \\
which for the love of Mankind He suffered.\evb\evg

\sectionline
